\documentclass[10pt,a4paper]{amsart}
\usepackage{amsmath,amssymb,amsthm,graphicx}
\newcommand{\R}{\mathbb{R}}
\newcommand{\Z}{\mathbb{Z}}
\newcommand{\N}{\mathbb{N}}
\newcommand{\C}{\mathbb{C}}
\newcommand{\E}{\mathcal{E}}
\DeclareMathOperator{\am}{am}
\DeclareMathOperator{\sn}{sn}
\DeclareMathOperator{\cn}{cn}
\DeclareMathOperator{\dn}{dn}

\numberwithin{equation}{section}

\newtheorem{Theorem}{Theorem}[section]
\newtheorem*{Theorem*}{Theorem}
\newtheorem{Corollary}[Theorem]{Corollary}
\newtheorem{Lemma}[Theorem]{Lemma}
\newtheorem{Proposition}[Theorem]{Proposition}
 { \theoremstyle{definition}
\newtheorem{Definition}[Theorem]{Definition}
\newtheorem{Note}[Theorem]{Note}
\newtheorem{Example}[Theorem]{Example}
\newtheorem{Remark}[Theorem]{Remark} }
\usepackage{cleveref}
\crefname{originallemma}{original lemma}{original lemmas}
\Crefname{originallemma}{Original lemma}{Original lemmas}
\makeatletter
\newlabel{l3}{{internal}{source}{}{sourcepointer-1}{}}
\newlabel{l3@cref}{{[originallemma][0][]internal}{[1][0][]source}{}{}{}}
\makeatother
\title{Continuity of the fundamental solutions in the eigenvalue parameter}
\author{Volkmer, Hans}
\date{doi:10.3842/SIGMA.2024.021; CC BY 4.0}
\begin{document}
\maketitle
\setcounter{section}{4}
\setcounter{Theorem}{0}
\noindent\textit{Editorial note.} All delivered mathematical text is the author's own, copied from the pinned SIGMA article source; no statement or proof step was rewritten, reordered or removed. The counted statement is the labeled theorem environment at source lines 319-327 and the counted proof is the proof environment at source lines 328-346.\par\medskip
\noindent\textit{Editorial note.} The retained preamble of this unit was generated mechanically for this delivery: it copies the author's own operator, macro, numbering and theorem declaration block from the pinned source (source lines 9-30), with the publisher class and page setup lines removed.\par\medskip
\noindent\textit{Allowed standard prerequisites (editorial).}\begin{itemize}
\item The representation of the solutions of the limiting equation by hypergeometric functions. same article, the representation introduced in the preceding section. Identifies the limit of the difference of the two solutions at infinity by evaluating the hypergeometric function at zero.
\item The explicit equation satisfied by the difference of the two solutions. derived at proof lines 337-338. Turns the problem into a forced harmonic oscillator equation whose solution is written by the variation of constants formula.
\item The uniform bound for the comparison solution. same article, the lemma at source lines 237-244. Bounds the forcing term and hence the integral that appears in the variation of constants formula.
\end{itemize}\medskip
\ifdefined\hypertarget\hypertarget{sourcepointer-1}{}\fi
\noindent\textit{Source reference.} Original lemma internal: sources/sigma-2024-021-on-the-hill-discriminant-of-lam-233-s-differ/pinned-original-source.tex, source lines 237--244. The proof invokes that lemma to bound the forcing term and the integral; the lemma is the author's own statement earlier in the same article and the estimate of the counted proof is retained verbatim.\par\medskip
\begin{Theorem}\label{t2}
Suppose $h>0$ and $h>\nu(\nu+1)$.
Then, for all $t\ge 0$,
\begin{gather*}
 |w_1(t)-z_1(t)|\le \omega^{-1}C_1|\nu|(\nu+1)(1-\tanh t),\\
 |w_2'(t)-z_2'(t)|\le C_2|\nu|(\nu+1)(1-\tanh t),
\end{gather*}
where $C_1$, $C_2$ are formed with $k=1$.
\end{Theorem}

\begin{proof}
Since $F(a,b;c;0)=1$, the representation $w_j(t)=\operatorname{Re}(v_j(t))$ yields
\begin{equation}\label{limit1}
\lim_{t\to\infty} w_j(t)-z_j(t)=\lim_{t\to\infty} \operatorname{Re}(A_j \bigl(2\cosh t)^{{\rm i}\omega}-A_j{\rm e}^{{\rm i}\omega t}\bigr) = 0.
\end{equation}
Similarly, we have
\begin{equation}\label{limit2}
 \lim_{t\to\infty} w_j'(t)-z_j'(t)=0 .
\end{equation}
The function $u_j=w_j-z_j$ satisfies
\[ u_j''+\omega^2 u_j=g_j(t), \qquad g_j(t):=\nu(\nu+1) \bigl(\tanh^2 t-1\bigr) w_j(t) .\]
Let $t_0, t\ge 0$. Then
\[ u_j(t)=u_j(t_0)\cos(\omega (t-t_0))+u_j'(t_0)\frac{\sin(\omega(t-t_0))}{\omega}+
\int_{t_0}^t \frac{\sin(\omega(t-s))}{\omega}g_j(s)\,{\rm d}s .\]
Letting $t_0\to\infty$, using \eqref{limit1}, \eqref{limit2} and Lemma \ref{l3}, we obtain
\[ |u_1(t)|\le\omega^{-1}\int_t^\infty |g_1(s)|\,{\rm d}s \le\omega^{-1}C_1|\nu|(\nu+1)(1-\tanh t) \]
as desired.
The estimate for $u_2'$ is derived similarly.
\end{proof}

\end{document}
